\documentclass[11pt, a4paper]{article}
\usepackage[a4paper, top=2.5cm, bottom=2.5cm, left=2cm, right=2cm]{geometry}
\usepackage{amsmath,amssymb,amsthm,microtype}
\usepackage{hyperref}

\theoremstyle{definition}
\newtheorem{problem}{Problem}
\theoremstyle{plain}
\newtheorem{theorem}{Theorem}
\newtheorem{lemma}{Lemma}

\begin{document}

\title{\vspace{-1.5cm}\textbf{Regional Quantitative Problem-Solving Circle} \\ \Large Solution Key 1: Extremal Principles \& Monovariants (Weeks 1--2)}
\author{\textbf{Lead Architect:} Mohamed Jad Srifi}
\date{\textbf{Term:} Fall 2025 -- Spring 2026}
\maketitle

\section*{Rigorous Proofs \& Structural Solutions}

\begin{problem}[Combinatorial Geometry / Extremal Distances]
$n \ge 3$ points in the plane have mutually distinct pairwise distances. Each point connects to its nearest neighbor. Prove there exists at least one pair of mutually connecting points, and prove that no edges intersect.
\end{problem}

\begin{proof}
\textbf{Part 1 (Mutual Connection):} 
Because the set of points $S$ is finite, the set of all pairwise distances is also finite. By the Extremal Principle, this finite set of positive real numbers must possess a strict global minimum. Let $A, B \in S$ be the two points separated by this absolute minimum distance $d(A, B)$. 
Because all distances are distinct, no other point $X \in S$ can be closer to $A$ than $B$ is, meaning $d(A, X) > d(A, B)$. Therefore, $A$'s nearest neighbor is strictly $B$, so $A$ connects to $B$. Symmetrically, no point is closer to $B$ than $A$ is, so $B$ connects to $A$. Thus, the pair $(A, B)$ mutually connects.

\textbf{Part 2 (No Intersections):}
Assume, for the sake of contradiction, that two segments $AB$ and $CD$ intersect at a point $X$. Without loss of generality, assume the edge $AB$ was created because $B$ is the nearest neighbor to $A$, and the edge $CD$ was created because $D$ is the nearest neighbor to $C$. 
By definition of the nearest neighbor connections, we must have $d(A, B) < d(A, D)$ and $d(C, D) < d(C, B)$.
Summing these two strict inequalities yields:
\[
d(A, B) + d(C, D) < d(A, D) + d(C, B)
\]
However, consider the geometry of the intersection at $X$. By the triangle inequality applied to $\triangle AXD$ and $\triangle CXB$:
\[
d(A, D) < d(A, X) + d(X, D) \quad \text{and} \quad d(C, B) < d(C, X) + d(X, B)
\]
Summing these inequalities gives:
\[
d(A, D) + d(C, B) < d(A, X) + d(X, D) + d(C, X) + d(X, B)
\]
Regrouping the terms on the right side using the collinearity of the segments ($d(A, X) + d(X, B) = d(A, B)$ and $d(C, X) + d(X, D) = d(C, D)$):
\[
d(A, D) + d(C, B) < d(A, B) + d(C, D)
\]
We have derived that $d(A, B) + d(C, D) < d(A, D) + d(C, B)$ and simultaneously $d(A, D) + d(C, B) < d(A, B) + d(C, D)$. This is a strict mathematical contradiction. Therefore, our assumption is false, and no two edges can intersect.
\end{proof}

\vspace{0.5cm}

\begin{problem}[The Sylvester-Gallai Theorem]
Let $S$ be a finite set of points in $\mathbb{R}^2$ not all collinear. Prove there exists an ordinary line passing through exactly two points using the minimum distance/area extremal argument.
\end{problem}

\begin{proof}
Assume, for the sake of contradiction, that no ordinary line exists. Therefore, every line that passes through at least two points of $S$ must pass through at least three points of $S$.

Consider the set of all pairs $(P, \ell)$, where $\ell$ is a line connecting at least two points of $S$, and $P \in S$ is a point not lying on $\ell$. Because $S$ is finite and the points are not all collinear, this set of pairs is non-empty.
By the Extremal Principle, there must exist a pair $(P_0, \ell_0)$ that strictly minimizes the perpendicular distance from the point to the line. Let this minimal positive distance be $d(P_0, \ell_0)$, with orthogonal projection $Q$ on $\ell_0$.

By our initial contradiction assumption, $\ell_0$ contains at least three points of $S$. By the Pigeonhole Principle on the 1D geometry of $\ell_0$, at least two of these points must lie on the same side of $Q$. Let these points be $A$ and $B$, ordered such that $A$ lies strictly between $Q$ and $B$.

Construct a new line $\ell_1$ passing through $P_0$ and $B$. Since $P_0, B \in S$, $\ell_1$ is a valid line generated by the set. Now consider the perpendicular distance from $A$ to $\ell_1$, denoted $d(A, \ell_1)$. 
In the right triangle $\triangle P_0 Q B$, the hypotenuse is $P_0 B$. The point $A$ lies on the leg $QB$. Basic geometry dictates that the altitude dropped from $A$ to the hypotenuse $\ell_1$ is strictly shorter than the altitude dropped from $P_0$ to $\ell_0$. 
Thus, $d(A, \ell_1) < d(P_0, \ell_0)$. However, $A \in S$ and $A \notin \ell_1$, meaning $(A, \ell_1)$ is a valid pair in our defined set with a strictly smaller distance than our assumed global minimum $(P_0, \ell_0)$. This is a contradiction. Thus, an ordinary line must exist.
\end{proof}

\vspace{0.5cm}

\begin{problem}[Infinite Lattice Harmonicity / Monovariant]
Each cell of an infinite square grid contains a positive integer equal to the arithmetic mean of its four orthogonal neighbors. Using the Well-Ordering Principle, prove all entries on the grid are identical.
\end{problem}

\begin{proof}
Let the infinite grid be defined by a function $v(x, y) \in \mathbb{Z}^+$ for coordinates $(x, y) \in \mathbb{Z}^2$.
Consider the set of all values present on the grid. Because this is a non-empty subset of the strictly positive integers, the Well-Ordering Principle guarantees that this set possesses a strict absolute minimum element. Let this minimal value be $m$.

There must exist at least one cell on the grid containing this minimal value. Let its coordinates be $(x_0, y_0)$, so $v(x_0, y_0) = m$.
By the harmonic property given in the problem, this value is the arithmetic mean of its four neighbors:
\[
m = \frac{v(x_0, y_0+1) + v(x_0, y_0-1) + v(x_0-1, y_0) + v(x_0+1, y_0)}{4}
\]
By our definition of $m$ being the absolute minimum, we know that $v(x_0, y_0+1) \ge m$, $v(x_0, y_0-1) \ge m$, $v(x_0-1, y_0) \ge m$, and $v(x_0+1, y_0) \ge m$.
If even a single one of these neighbors were strictly greater than $m$, the arithmetic mean of the four values would evaluate to a number strictly greater than $m$, breaking the equality. 

Therefore, to preserve the harmonic average $m$, it is algebraically forced that all four orthogonal neighbors must also exactly equal $m$. 
By mathematical induction propagating outward through orthogonal adjacencies, every connected cell in the infinite grid must share this value $m$. Thus, all integers on the grid are identical.
\end{proof}

\end{document}
